Predicted Signal for Goldstone Mode: If the anisotropy-gapped Goldstone mode of the spin cycloid is excited, the orientation of the ferroelectric polarization should oscillate. In principle, the magnitude of the polarization in its original orientation (along the y-axis) should decrease at a frequency twice as large as the oscillation, while an oscillating electric polarization should appear along the z-axis. It is unclear why such a signal is not observed in the S-in/S-out or P-in/P-out channels.
